% from ega4-15-terminal-fr.tex
\oldpage[IV-3]{244}
le point générique de $Y'$, alors, pour tout point $x'\in Z'\cap f'^{-1}(\eta')$ rationnel sur $\mathbf k(\eta')$, il existe une $Y'$-section de $X'$ contenant $x'$.
\end{enumerate}
\end{corollary}

\begin{proof}
Le fait que a) entraîne b) résulte de ce que $\overline Z_1\subset g_1^{-1}(\overline Z)$ et de ce que $g_1^{-1}(\overline Z)$ est propre sur $Y_1$ (15.7.1). Le fait que b) entraîne c) résulte de (\textbf{II}, 7.3.3). Reste donc à voir que c) implique a). Il suffit évidemment de prouver que toute composante irréductible de $\overline Z$ est propre sur $Y$ au point $y$, et l'hypothèse permet donc de se borner au cas où $Z$ est réduit à un seul point $z$. On peut évidemment aussi, compte tenu de (\textbf{I}, 5.2.2) et (\textbf{II}, 5.4.6), supposer que $X$ et $Y$ sont réduits, puis (par définition d'une partie propre (\textbf{II}, 5.4.10)) supposer que $\overline Z=X=\overline{\{z\}}$ et (appliquant encore (\textbf{I}, 5.2.2)) que $f$ est dominant, donc $X$ et $Y$ intègres et noethériens; il faut alors prouver que $f$ est propre au point $y$ de $Y$. Montrons pour cela que $f$ vérifie la condition c) de (15.7.2). Soit $A'$ un anneau de valuation discrète ayant $R(X)=\mathbf k(z)$ pour corps des fractions et dominant $\mathcal O_y$; avec les notations de c), on a donc $g(y')=y$, et $g(\eta')$ est le point générique $\eta=f(z)$ de $Y$. Comme par définition $\mathbf k(\eta')=\mathbf k(z)$, il existe un $\mathbf k(\eta)$-morphisme $\operatorname{Spec}(\mathbf k(\eta'))\to\operatorname{Spec}(\mathbf k(z))$ transformant $\eta'$ en $z$, donc une $\mathbf k(\eta')$-section de $f'^{-1}(\eta')$ (\textbf{I}, 3.3.14); si $z'$ est l'image de $\eta'$ par cette section, $z'$ est donc rationnel sur $\mathbf k(\eta')$ et $g'(z')=z$. On conclut donc de l'hypothèse c) qu'il existe une $Y'$-section $h'$ de $X'$ telle que $h'(\eta')=z'$; soit $h=g'\circ h':Y'\to X$ le $Y$-morphisme correspondant, et soit $x=h(y')$; il est clair que $f(x)=y$ et que $A'=\mathcal O_{y'}$ domine $\mathcal O_x$, ce qui achève la démonstration.
\end{proof}

\begin{corollary}[15.7.5]
\label{IV.15.7.5-fr}
Soient $Y$ un préschéma localement noethérien, $f:X\to Y$ un morphisme séparé de type fini, $y$ un point de $Y$. Les conditions suivantes sont équivalentes :
\begin{enumerate}
  \item[\rm a)] $f$ est propre au point $y$.
  \item[\rm b)] Si $Y_1=\operatorname{Spec}(\mathcal O_y)$, $X_1=X\times_Y Y_1$, le morphisme $f_1=f_{(Y_1)}:X_1\to Y_1$ est propre.
  \item[\rm c)] Pour tout schéma $Y'=\operatorname{Spec}(A')$, où $A'$ est un anneau de valuation discrète, et tout morphisme $g:Y'\to Y$, tel que l'image $g(y')$ du point fermé $y'$ de $Y'$ soit $y$, on a la propriété suivante : posant $X'=X\times_Y Y'$, $f'=f_{(Y')}:X'\to Y'$ et désignant par $\eta'$ le point générique de $Y'$, alors, pour tout $x'\in f'^{-1}(\eta')$, rationnel sur $\mathbf k(\eta')$, il existe une $Y'$-section de $X'$ contenant $x'$.
\end{enumerate}
\end{corollary}

\begin{corollary}[15.7.6]
\label{IV.15.7.6-fr}
Soient $Y$ un préschéma localement noethérien, $f:X\to Y$ un morphisme de type fini. Les conditions suivantes sont équivalentes :
\begin{enumerate}
  \item[\rm a)] Il existe un morphisme propre $g:X\to Z$ et une immersion $h:Z\to Y$ tels que $f=h\circ g$.
  \item[\rm b)] Le sous-espace $f(X)$ est localement fermé dans $Y$, et si $Z'$ est le sous-préschéma de $Y$ image fermée de $X$ par $f$ (\textbf{I}, 9.5.3 et 9.5.1), et $Z''$ le préschéma induit par $Z'$ sur l'ouvert $f(X)$ de $Z'$, de sorte que $f$ se factorise de façon unique en $f=j\circ g$ (où $j:Z''\to Y$ est l'injection canonique), alors $g:X\to Z''$ est propre.
  \item[\rm c)] Pour tout $y\in f(X)$, $f$ est propre au point $y$.
  \item[\rm d)] Pour tout schéma $Y'=\operatorname{Spec}(A')$, où $A'$ est un anneau de valuation discrète, et tout morphisme $g:Y'\to Y$ tel que $g(Y')\subset f(X)$, toute $Y'$-section rationnelle (\textbf{I}, 7.1.2) de $X'=X\times_Y Y'$ se prolonge de façon unique en une $Y'$-section de $X'$.
\end{enumerate}
\end{corollary}

\begin{proof}
Il est clair que b) entraîne a). Pour voir que a) entraîne c), remarquons d'abord que a) entraîne que $h(Z)$ est localement fermé dans $Y$ et $g(X)$ fermé dans $Z$, donc $f(X)$ est localement fermé dans $Y$; pour tout $y\in f(X)$, il y a par suite un voisinage ouvert $U$
